%% sections/02-probability-foundations.tex
%% Section 2: Probability Foundations (~7-9 pages)
%%
%% Build philosophy (from project owner): match the depth of v5 LaTeX notes,
%% not Joseph Genzer's shortened treatment. Build from the random variable
%% definition all the way to multivariate normal. Prove as much as possible.
%% Include figures wherever they help.
%%
%% Source notes:
%%   - v5 LaTeX notes: full MGF, full CF, modes of convergence, CLT via CFs,
%%     random vectors and PSD sections.
%%   - Handwritten Note 1.pdf, Note 2.pdf.
%%   - Casella and Berger (2002), Chapter 5; Joseph Genzer (2025), Section 2.

\section{Probability Foundations}\label{sec:foundations}

\todo{Opening paragraph: motivation. The high-dimensional phenomena we will
study in later sections (Marchenko-Pastur, Tracy-Widom, the spike phase
transition) all rest on a small set of foundational ideas from classical
probability. This section develops them in depth, starting from the
definition of a random variable. The depth here is deliberate: most of the
later analysis is, at its heart, a careful application of these basic tools
(differentiating under the integral, characteristic function expansions,
the multivariate central limit theorem, the spectral theorem). A reader who
is fluent with this material can skim; everyone else benefits from the
deliberate pace.}

\subsection{Random Variables and Distributions}\label{subsec:rvs}

\todo{Build absolute foundation (from v5 + Joseph Section 2.1):
\begin{itemize}
  \item Definition: probability space \((\Omega, \mathcal{F}, \Prob)\) at a light touch.
  \item Definition: random variable \(X: \Omega \to \R\).
  \item Definition: CDF \(F_X(x) = \Prob(X \le x)\) and its three characterizing properties (limits, right-continuity, monotonicity). State the converse (any function satisfying these is a CDF), cite \citep[][\S 1.2]{Casella2002}.
  \item Definition: discrete vs.\ absolutely continuous random variables; PMF and PDF.
  \item Examples: Bernoulli, binomial, Poisson, normal, gamma, Cauchy. Densities listed.
  \item Definition: independence of random variables.
\end{itemize}
Cite Casella-Berger and Joseph throughout.}

\subsection{Moments and Moment Generating Functions}\label{subsec:mgf}

\todo{Full depth from v5. Sub-structure:
\begin{itemize}
  \item Definition: \(j\)th moment, \(j\)th central moment, variance.
  \item Definition: MGF \(M_X(t) = \E[e^{tX}]\) with existence requirement on an open neighborhood of zero.
  \item Heuristic: MGF as the two-sided Laplace transform of the density (use this to motivate uniqueness).
  \item Theorem: if \(M_X\) exists in a neighborhood of \(0\), then \(M_X^{(n)}(0) = \E[X^n]\). Full proof by differentiation under the integral (continuous case) and term-by-term differentiation (discrete case).
  \item Theorem: MGF uniqueness. If \(M_X(t) = M_Y(t)\) on a neighborhood of \(0\), then \(X \eqd Y\). State precisely; sketch the easy direction; cite \citep[Thm.~2.3.11]{Casella2002} for the converse.
  \item Worked examples: MGFs of binomial, Poisson, normal, gamma. Derive each. Tabulate the resulting means and variances.
  \item Cautionary example: Cauchy has no MGF (the integral diverges for all \(t \ne 0\)). Use this to motivate characteristic functions in \S\ref{subsec:cf}.
\end{itemize}}

\subsection{Characteristic Functions}\label{subsec:cf}

\todo{Full depth from v5. Sub-structure:
\begin{itemize}
  \item Definition: \(\varphi_X(t) = \E[e^{i t X}]\). Always exists since \(\abs{e^{itX}} = 1\); deduce \(\abs{\varphi_X(t)} \le 1\) and \(\varphi_X(0) = 1\).
  \item Identification as a Fourier transform of the density (when one exists), and as a Fourier-Stieltjes transform of the distribution in general.
  \item Relation to the MGF when both exist: \(\varphi_X(t) = M_X(i t)\).
  \item Worked examples: derive \(\varphi_X\) for the standard normal (completing the square gives \(e^{-t^2/2}\)) and the Cauchy (gives \(e^{-|t|}\)). The Cauchy case is the punchline: characteristic functions handle heavy tails that defeat MGFs.
  \item Properties: independence implies \(\varphi_{X+Y}(t) = \varphi_X(t) \varphi_Y(t)\); linear transformations \(\varphi_{aX+b}(t) = e^{itb} \varphi_X(at)\).
  \item Statement of the uniqueness theorem for CFs (the Fourier inversion content).
  \item Statement of L\'evy's continuity theorem: pointwise convergence of CFs to a function continuous at \(0\) is equivalent to convergence in distribution. Defer the full proof to Appendix~\ref{app:auxiliary}; cite \citep[Thm.~5.5.15]{Casella2002}.
\end{itemize}
Suggested figure: side-by-side plot of \(\varphi_{\mathcal{N}(0,1)}(t) = e^{-t^2/2}\) and \(\varphi_{\mathrm{Cauchy}}(t) = e^{-|t|}\), to make tangible the contrast between light-tailed and heavy-tailed behavior at the CF level.}

\subsection{Modes of Convergence}\label{subsec:convergence}

\todo{Content (from v5 + Joseph Section 2.2):
\begin{itemize}
  \item Definitions:
    \(X_n \convas X\) (almost-sure convergence),
    \(X_n \convp X\) (convergence in probability),
    \(X_n \convd X\) (convergence in distribution).
  \item Theorem: implication chain a.s. \(\Rightarrow\) probability \(\Rightarrow\) distribution. Prove the first implication via the Borel-Cantelli lemma; sketch the second.
  \item Counterexamples to the converses:
    \begin{itemize}
      \item \(X_n \in \{0, 1\}\) with \(\Prob(X_n = 1) = 1/n\) converges in probability to \(0\) but not almost surely (the second Borel-Cantelli lemma forces \(X_n = 1\) infinitely often).
      \item For convergence in distribution but not in probability, use \(X \sim \mathcal{N}(0,1)\) and \(X_n = -X\) (the distribution is symmetric).
    \end{itemize}
  \item Continuous mapping theorem: if \(X_n \convas X\) and \(g\) is continuous, then \(g(X_n) \convas g(X)\); same for convergence in probability. Defer Slutsky's theorem to Appendix~\ref{app:auxiliary}.
\end{itemize}}

\subsection{Limit Theorems: LLN and CLT}\label{subsec:limit-theorems}

\todo{Content (from v5 + Joseph Section 2.2):
\begin{itemize}
  \item Inequalities first: Markov, then Chebyshev (the latter as a Markov application to \((X-\mu)^2\)). Full proofs.
  \item Chebyshev's weak law of large numbers: iid with finite variance \(\Rightarrow\) \(\bar X_n \convp \mu\). Full proof via Chebyshev.
  \item Statement of Kolmogorov's strong law of large numbers; cite \citep{Billingsley1995} for the proof.
  \item Central limit theorem: for iid with finite variance,
    \(\sqrt{n}(\bar X_n - \mu)/\sigma \convd \mathcal{N}(0, 1)\). Prove via characteristic function expansion using the L\'evy continuity theorem from \S\ref{subsec:cf}. Carry out the \((1 - t^2/(2n) + o(1/n))^n \to e^{-t^2/2}\) step explicitly.
\end{itemize}
Suggested figure (notebook 01): the WLLN convergence for normal and exponential samples, plotting \(\bar X_n\) versus \(n\). Reproduce Joseph's Figure~1 in our color palette. Also a second figure showing the CLT convergence: standardized sample means at \(n = 10, 100, 1000\) versus the standard normal density.}

\subsection{Random Vectors and the Multivariate Normal Distribution}\label{subsec:rvecs}

\todo{Content (from v5 + Joseph Section 2.4):
\begin{itemize}
  \item Definition: random vector \(\vect X = (X_1, \dots, X_p)\transpose\), joint CDF, joint PMF / PDF.
  \item Mean vector \(\vect \mu = \E[\vect X]\) and covariance matrix \(\Sigma = \E[(\vect X - \vect \mu)(\vect X - \vect \mu)\transpose]\).
  \item Theorem: \(\Sigma\) is positive semidefinite. Full proof via \(v\transpose \Sigma v = \Var(v\transpose \vect X) \ge 0\).
  \item When is \(\Sigma\) positive definite vs.\ semidefinite? Connect to the dimension of the support.
  \item Definition: multivariate normal distribution via density. State the equivalent characterization via characteristic function \(\varphi_{\vect X}(\vect t) = \exp(i \vect t\transpose \vect \mu - \tfrac12 \vect t\transpose \Sigma \vect t)\); this is the cleanest definition because it handles singular \(\Sigma\) without modification.
  \item Properties: linear transformations of multivariate normals are multivariate normal; uncorrelated jointly Gaussian variables are independent (special fact for the Gaussian).
  \item Statement of the multivariate central limit theorem; the proof is the same Cram\'er-Wold device applied to the univariate CLT, defer to Appendix~\ref{app:auxiliary}.
\end{itemize}}

\clearpage
